\section*{\texorpdfstring{\S\CurrentExerciseGroup}{Section \CurrentExerciseGroup}}
\phantomsection
\addcontentsline{toc}{section}{\texorpdfstring{Exercises for \S\CurrentExerciseGroup}{Exercises for Section \CurrentExerciseGroup}}

\begin{enumerate}
\def\labelenumi{\arabic{enumi})}
\item
  Let E be the vector space of functions \(x \mapsto g(x) = \int_{0}^{x} f(t) dt\) defined on the interval \([0,1]\) of R, where f runs over the set of regulated functions on \([0,1]\) . Let P be the set of increasing functions belonging to E. Show that P is a convex cone such that E = P - P, \(P \cap (-P) = \{0\}\) , and that E, equipped with the order structure defined by the relation \(g - h \in P\) , is a Riesz space.
\item
  Let \(I \subset \mathbf{R}\) be a compact interval. Show that the subspace of \(\mathbf{R}^I\) formed by the restrictions to \(I\) of the polynomial functions (with real coefficients) is a directed ordered space but is not a Riesz space.
\item
  \begin{enumerate}
  \def\labelenumii{\alph{enumii})}
  \tightlist
  \item
    Let E be a Riesz space, H a linear subspace of E. If, for any pair of elements x, y of H, \(\sup(x,y)\) belongs to H, then H is said to be a co-lattice subspace. For this to be the case, it is necessary and sufficient that the relation \(x \in H\) imply \(|x| \in H\) .
  \end{enumerate}
\end{enumerate}

\begin{enumerate}
\def\labelenumi{\alph{enumi})}
\setcounter{enumi}{1}
\tightlist
\item
  Let I be a compact interval in R, H the subspace of \(R^{I}\) formed by the restrictions to I of the first-degree polynomials \(t \mapsto \alpha t + \beta\) . Show that H is not a co-lattice subspace of \(R^{I}\) but that H is a Riesz space (for the structure induced by that of \(R^{I}\) ).
\end{enumerate}

\begin{enumerate}
\def\labelenumi{\arabic{enumi})}
\setcounter{enumi}{3}
\item
  Let E be a Riesz space, H a linear subspace of E. One says that H is an isolated subspace of E if the relations \(x \in H\) , \(|y| \leqslant |x|\) imply \(y \in H\) . In this case, let P be the set of elements \(\dot{x}\) of the quotient space E/H such that there exists at least one element \(x \geqslant 0\) in the class \(\dot{x}\) . Show that P is the set of elements \(\geqslant 0\) for an order structure on E/H, compatible with the vector space structure of E/H, for which E/H is a Riesz space (cf.~A, VI, §1, Exer. 4).
\item
  \begin{enumerate}
  \def\labelenumii{\alph{enumii})}
  \tightlist
  \item
    An ordered vector space E is said to be archimedean if every \(x \in E\) such that the set of nx (n an integer \(\geqslant 0\) ) is bounded above, is necessarily \(\leqslant 0\) (A, VI, §1, Exer. 31). Show that this condition is equivalent to the following: the intersection of any plane passing through 0 with the convex cone P of elements \(\geqslant 0\) of E is a closed angular sector. Show that for an ordered vector space to be archimedean, it is necessary and sufficient that it be isomorphic to a subspace of a fully lattice-ordered space (cf.~A, VI, §1, Exer. 31).
  \end{enumerate}
\end{enumerate}

\begin{enumerate}
\def\labelenumi{\alph{enumi})}
\setcounter{enumi}{1}
\tightlist
\item
  Let F be the Riesz space \(\mathbf{R}^{\mathbf{R}}\) of real-valued functions defined on \(\mathbf{R}\) and let H be the subspace \(\mathcal{B}(\mathbf{R})\) of bounded functions on \(\mathbf{R}\) ; H is an isolated subspace of F (Exer. 4). Show that the Riesz space \(\mathrm{E} = \mathrm{F} / \mathrm{H}\) (Exer. 4) is not archimedean: more precisely, show that for every element \(x \geqslant 0\) of E, there exists a \(y \in \mathrm{E}\) such that \(y \geqslant nx\) for every integer \(n \geqslant 0\) .
\end{enumerate}

\begin{enumerate}
\def\labelenumi{\arabic{enumi})}
\setcounter{enumi}{5}
\item
  Let E be a fully lattice-ordered space, V a linear subspace of E. Show that, for there to exist in E a supplement W of V such that E is the ordered direct sum of V and W, it is necessary and sufficient that V be a band; W is then necessarily identical to the band of elements of E that are alien to every element of V.
\item
  Let E be a fully lattice-ordered space, H an isolated subspace of E (Exer. 4). If \(B_{1}\) and \(B_{2}\) are two supplementary bands in E, show that H is the ordered direct sum of \(H_{1}=H\cap B_{1}\) and \(H_{2}=H\cap B_{2}\) , and that the Riesz space E/H is the ordered direct sum of the Riesz spaces \(B_{1}/H_{1}\) and \(B_{2}/H_{2}\) .
\item
  Let \((\mathbf{E}_{\iota})\) be a family of fully lattice-ordered spaces, E the product space of the \(E_{\iota}\) (which is fully lattice-ordered). Show that if B is a band in E, then each of its projections \(\mathrm{B}_{\iota} = \mathrm{pr}_{\iota}(\mathrm{B})\) is a band in \(E_{\iota}\) , and B is identical to the product of the \(B_{\iota}\) . From this, deduce the determination of the bands in the space \(R^{A}\) of mappings of a set A into R. Show that, in the space \(\mathcal{B}(A)\) of bounded real-valued functions on A, every band is the trace on \(\mathcal{B}(A)\) of a band in \(R^{A}\) .
\item
  Let E be a fully lattice-ordered space. A filter \(\mathfrak{F}\) on E is said to be bounded above (resp. bounded below, bounded) if there is a set in \(\mathfrak{F}\) that is bounded above (resp. bounded below, bounded). For a filter \(\mathfrak{F}\) that is bounded above (resp. below), the limit superior (resp. limit inferior) of \(\mathfrak{F}\) , denoted lim sup \(\mathfrak{F}\) (resp. lim inf \(\mathfrak{F}\) ) is defined to be the element \(\inf_{\mathrm{X}}(\sup \mathrm{X})\) (resp. \(\sup_{\mathrm{X}}(\inf \mathrm{X})\) ) of E, where X runs over the set of sets
\end{enumerate}

in \(\mathfrak{F}\) that are bounded above (resp. bounded below). If \(\mathfrak{F}\) is bounded, then \(\liminf\mathfrak{F} \leqslant \limsup\mathfrak{F}\) ; \(\mathfrak{F}\) is said to have a limit, or to be convergent, for the order structure of E, if \(\limsup\mathfrak{F} = \liminf\mathfrak{F}\) ; the common value of these two elements is then denoted \(\lim\mathfrak{F}\) , and \(\mathfrak{F}\) is said to converge to this element of E. If a bounded filter has a limit, then every finer filter has the same limit (for the order structure).

\begin{enumerate}
\def\labelenumi{\alph{enumi})}
\tightlist
\item
  In order that a bounded filter \(\mathfrak{F}\) have a limit for the order structure, it is necessary and sufficient that \(\inf_{\mathbf{X}} (\sup \mathbf{X} - \inf \mathbf{X}) = 0\) as \(\mathbf{X}\) runs over the set of bounded elements of \(\mathfrak{F}\) (one can make use of the fact that if, in E, \((x_{\alpha})\) and \((y_{\alpha})\) are two decreasing families that are bounded below, having the same right-directed index set, then
\end{enumerate}

\[
\inf (x _ {\alpha} + y _ {\alpha}) = \inf x _ {\alpha} + \inf y _ {\alpha};
\]

to prove this formula, one observes that \(\inf (x_{\alpha} + y_{\alpha})\leqslant x_{\beta} + y_{\gamma}\) for every pair of indices \(\beta ,\gamma)\) .

\begin{enumerate}
\def\labelenumi{\alph{enumi})}
\setcounter{enumi}{1}
\item
  Let A be a set filtered by a filter G, f a mapping of A into E; f is said to have a limit with respect to G for the order structure of E if \(f(\mathfrak{G})\) is a base of a bounded filter having a limit in E (for the order structure). Show that if A is a directed ordered set and f is an increasing mapping of A into E that is bounded above on A, then f has a limit with respect to the section filter of A, equal to \(\sup f(x)\) .
\item
  Let \((\mathbf{E}_{\iota})\) be a family of fully lattice-ordered spaces, E the product fully lattice-ordered space of the \(E_{\iota}\) . In order that a bounded filter F on E be convergent for the order structure, it is necessary and sufficient that, for every \(\iota\) , the filter with base \(\Pr_{\iota}(\mathfrak{F})\) be convergent in \(E_{\iota}\) ; if \(a_{\iota}\) is its limit, then \(a = (a_{\iota})\) is the limit of F.
\end{enumerate}

¶ 10) a) Let E be a fully lattice-ordered space; there exists a topology \(\mathcal{T}_0(E)\) on E that is the finest of the topologies \(\mathcal{T}\) on E for which every bounded filter \(\mathfrak{F}\) on E that converges in the sense of the order structure (Exer. 9) converges to the same limit for \(\mathcal{T}\) .

Let E and F be two fully lattice-ordered spaces, f a mapping of E into F such that, for every bounded filter F on E that is convergent for the order structure, \(f(\mathfrak{F})\) is a base of a bounded filter on F that converges to \(f(\lim \mathfrak{F})\) for the order structure. Show that, under these conditions, f is continuous for the topologies \(\mathcal{T}_{0}(\mathrm{E})\) and \(\mathcal{T}_{0}(\mathrm{F})\) .

\begin{enumerate}
\def\labelenumi{\alph{enumi})}
\setcounter{enumi}{1}
\item
  Deduce from a) that the topology \(\mathcal{T}_{0}(\mathrm{E})\) is compatible with the vector space structure of E and that the mapping \(x \mapsto |x|\) is continuous for this topology. Prove that \(\mathcal{T}_{0}(\mathrm{E})\) is compatible with the ordered vector space structure of E and, from this, deduce that this topology is Hausdorff.
\item
  Given any infinite set A, let \(\mathrm{E} = \mathcal{B}(\mathrm{A})\) be the fully lattice-ordered space of bounded real-valued functions on A. Let \(\Phi\) be the set of numerical functions \(\varphi\) (finite or not) defined on A such that \(\varphi(t) > 0\) for all \(t \in A\) and such that, for every integer n \textgreater{} 0, the set of \(t \in A\) such that \(\varphi(t) \leqslant n\) is finite. For every function \(\varphi \in \Phi\) , let \(V_{\varphi}\) be the set of \(x \in E\) such that \(|x| \leqslant \varphi\) ; show that the sets \(V_{\varphi}\) form a fundamental system of neighborhoods of 0 for a topology \(\mathcal{T}_{1}(\mathrm{E})\) compatible with the ordered vector space structure of E, and for which E is Hausdorff and complete. Show that the topology \(\mathcal{T}_{0}(\mathrm{E})\) is finer than \(\mathcal{T}_{1}(\mathrm{E})\) , but strictly coarser than the topology of uniform convergence in A (defined by the norm \(\|x\| = \sup_{t \in A} |x(t)|\) ) (consider the elements \(1 - \varphi_{X}\) of \(\mathcal{B}(\mathrm{A})\) ,
\end{enumerate}

where X runs over the set of finite subsets of A, \(\varphi_{X}\) being the characteristic function of X). Deduce from this that in order for a subset of E to be bounded (for the order structure), it is necessary and sufficient that it be bounded for the topology \(\mathcal{T}_{0}(E)\) (TVS, III, §1, No.~2). Show that every filter on E that is bounded and convergent for the topology \(\mathcal{T}_{0}(E)\) is convergent to the same limit for the order structure. Finally, show that there exist filters on E that are convergent for \(\mathcal{T}_{0}(E)\) but are not bounded.

¶ 11) Let E be a fully lattice-ordered space, \((x_{\iota})_{\iota \in I}\) a family of elements of E. For every finite subset H of I, set \(s_{H} = \sum_{\iota \in H} x_{\iota}\) ; the family \((x_{\iota})_{\iota \in I}\) is said to be summable

for the order structure of E if the mapping \(\mathrm{H} \mapsto s_{\mathrm{H}}\) has a limit for this order structure, with respect to the directed ordered set \(\mathfrak{F}(\mathrm{I})\) of finite subsets of I; this limit \(s\) is then called the sum of the family \((x_{\iota})_{\iota \in \mathrm{I}}\) and is denoted \(\sum_{\iota \in \mathrm{I}} x_{\iota}\) .

\begin{enumerate}
\def\labelenumi{\alph{enumi})}
\item
  For a family \((x_{\iota})_{\iota \in I}\) to be summable, it is necessary and sufficient that: \(1^{\circ}\) for every finite subset H of I, the set of \(|s_{K}|\) , where K runs over the set of finite subsets of I that do not intersect H, has a supremum \(r_{\mathrm{H}}\) in E; \(2^{\circ} \inf_{\mathrm{H}} r_{\mathrm{H}} = 0\) as H runs over \(\mathfrak{F}(I)\) (make use of Exer. 9 a)).
\item
  Generalize to summable families, for the order structure of E, the properties of summable families in commutative topological groups (GT, III, §5, Props. 2, 3 and Th. 2).
\item
  Let \((x_{\iota})_{\iota \in I}\) be a family of elements \(\geqslant 0\) of E; for the family to be summable, it is necessary and sufficient that the set of finite partial sums \(s_{\mathrm{H}}\) be bounded above (make use of Exer. 9 b)). If \((x_{\iota})_{\iota \in I}\) is summable and if \((y_{\iota})_{\iota \in I}\) is a family of elements such that \(0 \leqslant y_{\iota} \leqslant x_{\iota}\) for all \(\iota\) , show that the family \((y_{\iota})\) is summable and that \(\sum_{\iota \in I} y_{\iota} \leqslant \sum_{\iota \in I} x_{\iota}\) , with equality holding only if \(x_{\iota} = y_{\iota}\) for all \(\iota \in I\) .
\item
  Let \((x_{\iota})\) be a family of elements of E; show that if the family \((|x_{\iota}|)\) is summable for the order structure, then so is the family \((x_{\iota})\) .
\end{enumerate}

\begin{enumerate}
\def\labelenumi{\arabic{enumi})}
\setcounter{enumi}{11}
\tightlist
\item
  Let \(\mathbf{E}\) be a fully lattice-ordered space. Show that there exists a family \((u_{\iota})\) of elements \(> 0\) of \(\mathbf{E}\) such that \(u_{\iota}\) and \(u_{\kappa}\) are alien for every pair of distinct indices \(\iota, \kappa\) , and such that for every \(x > 0\) in \(\mathbf{E}\) there exists at least one index \(\iota\) such that \(\inf(x, u_{\iota}) > 0\) (use Zorn's lemma). From this, deduce that for every \(x \geqslant 0\) there exists one and only one family \((x_{\iota})\) of elements \(\geqslant 0\) of \(\mathbf{E}\) such that, for every \(\iota, x_{\iota}\) belongs to the band \(B_{\iota}\) generated by \(u_{\iota}\) , and such that \(x = \sum_{\iota} x_{\iota}\) for the order structure of \(\mathbf{E}\) .
\end{enumerate}

(take \(x_{\iota}\) to be the component of \(x\) in the band \(B_{\iota}\) ). Conversely, every family \((x_{\iota})\) of elements \(\geqslant 0\) of E that is bounded above, such that \(x_{\iota} \in B_{\iota}\) for all \(\iota\) , is summable in E (Exer. 11).

¶ 13 a) Let E be a fully lattice-ordered space. If \(u \neq 0\) is an arbitrary element of E, show that the set of \(x \in E\) such that there exists an integer \(n\) for which \(|x| \leqslant n|u|\) , is a linear subspace \(C_u\) that is a fully lattice-ordered and co-lattice subspace (Exer. 3) of E. For every \(x \in C_u\) , one denotes by \(\|x\|\) the infimum of the scalars \(\lambda > 0\) such that \(|x| \leqslant \lambda |u|\) ; show that \(\|x\|\) is a norm on \(C_u\) (make use of the fact that E is archimedean (Exer. 5)), that \(C_u\) , equipped with this norm, is complete, and that \(\|x\| = \sup(\|x^+\|, \|x^-\|)\) .

\begin{enumerate}
\def\labelenumi{\alph{enumi})}
\setcounter{enumi}{1}
\item
  Let \(\mathcal{I}\) be the set of components of \(u\) in the bands of the space \(\mathrm{C}_u\) ; show that \(\mathcal{I}\) is the set of elements \(c\) of \(\mathrm{C}_u\) such that \(0 \leqslant c \leqslant u\) and \(\inf(c, u - c) = 0\) (use Riesz's theorem). Show that \(\mathcal{I}\) is closed in the normed space \(\mathrm{C}_u\) and that, for the order relation induced by that of \(\mathrm{C}_u\) , \(\mathcal{I}\) is a complete Boolean algebra (S, III, §1, Exers. 11 and 17; it suffices to show that if \((c_\iota)\) is a family of elements of \(\mathcal{I}\) then \(\sup_{\iota} c_\iota\) belongs to \(\mathcal{I}\) ).
\item
  Let \(x\) be any element of \(\mathbf{C}_u\) ; for every \(\lambda \in \mathbf{R}\) let \(c(\lambda)\) be the component of \(u\) in the band of \(\mathbf{C}_u\) generated by \((\lambda u - x)^+\) . Show that if \(\lambda \leqslant \mu\) then \(c(\lambda) \leqslant c(\mu)\) ; \(c(\lambda) = 0\) for \(\lambda < -\|x\|\) , and \(c(\lambda) = u\) for \(\lambda > \|x\|\) . Show that if \(c \in \mathcal{I}\) is such that \(c \leqslant c(\lambda)\) , then the component of \(x\) in the band generated by \(c\) is \(\leqslant \lambda c\) (observe that the bands generated by \(c(\lambda)\) and by \((\lambda u - x)^+\) are identical, and that the component of \(\lambda u - x\) in this band is equal to that of \((\lambda u - x)^+\) ; from this, deduce that the component of \(x\) in this same band is \(\leqslant \lambda c(\lambda)\) ). Similarly, show that if \(c \in \mathcal{I}\) is such that \(c \leqslant u - c(\lambda)\) , then the component of \(x\) in the band generated by \(c\) is \(\geqslant \lambda c\) .
\item
  It is known (GT, II, §4, Exer. 12) that there exists an order structure isomorphism \(c \mapsto \theta_c\) of the Boolean algebra \(\mathcal{I}\) onto the Boolean algebra formed by the characteristic functions of the sets in a compact, totally disconnected space S that are both open and closed. Show that the relation \(\sum_{i} \lambda_i c_i = 0\) implies that the function \(\sum_{i} \lambda_i \theta_{c_i}\) is 0 on S (make use of the decomposition lemma); deduce from this remark and c) that the mapping \(c \mapsto \theta_{c}\) may be extended to an isomorphism \(x \mapsto \theta_{x}\) of the normed space \(C_{u}\) onto the normed space \(\mathcal{C}(\mathrm{S};\mathbf{R})\) of continuous real-valued functions on S, so that \((\theta_{x})^{+} = \theta_{x^{+}}\) . (With the help of c), show that for every \(x \in C_{u}\) and every \(\varepsilon > 0\) , there exists an increasing sequence \((\lambda_{i})_{0 \leqslant i \leqslant n}\) of real numbers such that \(\lambda_{i} - \lambda_{i-1} \leqslant \varepsilon\)
\end{enumerate}

and \(0 \leqslant x - \sum_{i=1}^{n} \lambda_{i-1} (c(\lambda_{i-1}) - c(\lambda_i)) \leqslant \varepsilon u\) .

\begin{enumerate}
\def\labelenumi{\alph{enumi})}
\setcounter{enumi}{4}
\item
  In order that S be finite, it is necessary and sufficient that \(C_{u}\) be finite-dimensional; from this, deduce (with the help of Exer. 12) that every finite-dimensional fully lattice-ordered space is isomorphic to a product space \(R^{n}\) (cf.~§2, Exer. 7).
\item
  Show that S is an extremely disconnected space (GT, I, §11, Exer. 21) (make use of the fact that \(\mathcal{I}\) is a complete Boolean algebra). An extremely disconnected compact space is called a Stone space. Show that if X is a Stone space then, for every lower semi-continuous numerical function \(f\) on X, the upper semi-continuous regularization \(g\) of \(f\) (GT, IV, §6, No.~2) is continuous (for every \(a < g(x)\) show that \(x\) cannot belong to the closure of the (open) set of \(y\) such that \(g(y) < a\) , by observing that \(f(z) \leqslant a\) at a point \(z\) in the closure of this set). From this, deduce that \(\mathcal{C}(\mathrm{X};\mathbf{R})\) is then a fully lattice-ordered space. Show that when the Stone space X is infinite, the supremum in \(\mathcal{C}(\mathrm{X};\mathbf{R})\) of a set that is bounded above is not necessarily equal to its upper envelope (consider a non-closed open set in X); however, these two functions are equal on the complement of a meager set (consider the set of points where their difference is \(\geqslant 1/n\) ).
\item
  Let X be a Stone space; show that if f is a bounded real-valued function, defined on the complement of a nowhere dense subset M of X and continuous on X - M, then \(f\) may be extended by continuity to all of X (make use of the fact that \(\mathcal{C}(X; \mathbf{R})\) is fully lattice-ordered). From this, deduce that the set \(\mathcal{Q}(X; \overline{\mathbf{R}})\) , of continuous numerical functions \(f\) on X (finite or not) such that \(f(+\infty)\) and \(f(-\infty)\) are nowhere dense in X, is a fully lattice-ordered vector space.
\item
  Show that the band \(B_{u}\) generated by u in E is isomorphic (as an ordered space) to an isolated subspace (Exer. 4) of \(\mathcal{Q}(\mathrm{X};\overline{\mathbf{R}})\) (regard an element \(f \geqslant 0\) of \(\mathcal{Q}(\mathrm{X};\overline{\mathbf{R}})\) as the supremum of the elements \(\inf(f,n)\) ).
\end{enumerate}

\begin{enumerate}
\def\labelenumi{\arabic{enumi})}
\setcounter{enumi}{13}
\tightlist
\item
  Let E be a Riesz space, equipped with a Hausdorff topology compatible with the ordered vector space structure of E.
\end{enumerate}

\begin{enumerate}
\def\labelenumi{\alph{enumi})}
\item
  Let K be a compact subset of E, H a subset of K that is directed for the relation \(\leqslant\) ; show that the section filter F of H is convergent. (First prove that the set L of upper bounds of H belonging to K is nonempty and compact, then that the set of cluster points of F is contained in L, and finally that there cannot exist two distinct cluster points.)
\item
  Deduce from \(a\) ) that if, in \(\mathbf{E}\) , every interval \([a,b]\) is compact, then \(\mathbf{E}\) is fully lattice-ordered.
\end{enumerate}

\setcounter{exercisegroup}{2}
\edef\CurrentExerciseGroup{\arabic{exercisegroup}}
